% =====================================================================
%  Основное содержимое отчёта.
% =====================================================================

\labpart{Краткое описание MongoDB}

\texttt{MongoDB} -- документоориентированная СУБД: данные хранятся не в
таблицах со строками, а в коллекциях JSON-подобных документов (BSON).
В лабораторной работе используется официальный Docker-образ
\texttt{mongo:7}. По умолчанию свежий экземпляр MongoDB принимает любые
подключения без аутентификации -- контроль доступа необходимо включить
явно ключом \texttt{--auth} и настроить отдельно от установки. Модель прав
в MongoDB строится вокруг ролей: встроенные роли (например,
\texttt{root}, \texttt{readWrite}) назначаются пользователю в контексте
конкретной базы данных, а для более точных ограничений можно определить
собственную роль с перечнем разрешённых действий над конкретной
коллекцией -- этим МongoDB напоминает списки контроля доступа из первой
части лабораторной работы, только уровнем выше: не файл и не каталог, а
документ и коллекция.

\labpart{Краткое описание OpenFire}

\texttt{OpenFire} -- сервер обмена мгновенными сообщениями по протоколу
XMPP (Jabber), написанный на Java. Он состоит из двух независимых
подсистем контроля доступа: во-первых, веб-консоль администрирования
(порт 9090), доступ к которой разрешён только учётным записям из списка
администраторов; во-вторых, сам протокол XMPP, в котором права
участников многопользовательских чат-комнат (Multi-User Chat, MUC)
определяются их \emph{аффилиацией} (владелец, администратор, участник,
не участвует) и производной от неё \emph{ролью} в конкретной комнате
(модератор, участник, посетитель). В комнате с включённой модерацией
только участники с ролью не ниже «участник» имеют право голоса --
посетители (в том числе анонимные) могут только читать переписку.

\labpartbreak{Реализация: MongoDB}

Проект \texttt{Lab1.2-MongoDB} организован так же, как и программа из
первой части лабораторной работы: \texttt{Program.cs} последовательно
вызывает четыре этапа из каталога \texttt{Steps} (\texttt{Install},
\texttt{Configure}, \texttt{Verify}, \texttt{Cleanup}), а вспомогательный
класс \texttt{Report} запускает внешние процессы и печатает результат
каждой команды в формате \texttt{[OK]/[FAIL]}, а результат проверки прав
доступа -- в формате \texttt{[ALLOWED]/[DENIED]}.

\subsection*{1. Установка}

Установка выполняется командами \texttt{docker pull mongo:7} и
\texttt{docker run}, с монтированием именованного тома
\texttt{mongo\_lab\_data} для постоянного хранения данных. Контейнер
намеренно запускается без \texttt{--auth}: MongoDB разрешает создать
самого первого пользователя без пароля только через т.\,н. localhost
exception, то есть до включения обязательной аутентификации.

\subsection*{2. Настройка контроля доступа}

Пока аутентификация ещё не включена, программа создаёт администратора
\texttt{admin} с встроенной ролью \texttt{root}, после чего контейнер
перезапускается с ключом \texttt{--auth} (том с данными переиспользуется,
поэтому созданный пользователь сохраняется). Далее от имени
администратора создаются: пользователь \texttt{labuser} с ролью
\texttt{readWrite} на базе \texttt{labdb} (полные права в предоставленной
администратором «подсистеме» -- одной прикладной базе данных) и
пользователь \texttt{labguest}, которому назначена не встроенная, а
специально созданная роль \texttt{guestReader}, разрешающая только
действие \texttt{find} и только над одной коллекцией
\texttt{public\_notes} (см. \autoref{lst:mongo-guest-role}). Соседняя
коллекция \texttt{private\_notes} для гостя недоступна вовсе.

\begin{codelisting}{csharp}{Создание гостевой роли и пользователя с правом чтения одной коллекции}{mongo-guest-role}
Report.Step(8, $"Create '{Config.GuestUser}' — read-only on '{Config.AppDb}.{Config.PublicCollection}' only");
MongoEvalAuth($"Create role '{Config.GuestRole}' (find on {Config.AppDb}.{Config.PublicCollection} only)",
    $"db.getSiblingDB('{Config.AppDb}').createRole({{role:'{Config.GuestRole}'," +
    $"privileges:[{{resource:{{db:'{Config.AppDb}',collection:'{Config.PublicCollection}'}},actions:['find']}}],roles:[]}})");
MongoEvalAuth($"Create user '{Config.GuestUser}' (role {Config.GuestRole})",
    $"db.getSiblingDB('{Config.AppDb}').createUser({{user:'{Config.GuestUser}',pwd:'{Config.GuestPassword}',roles:[{{role:'{Config.GuestRole}',db:'{Config.AppDb}'}}]}})");
\end{codelisting}

\subsection*{3. Проверка корректности}

Проверка выполняется той же утилитой \texttt{mongosh}, но уже от имени
каждой из трёх учётных записей, и интерпретирует нулевой код завершения
команды как \texttt{ALLOWED}, а ненулевой (сервер отклонил операцию как
неавторизованную) -- как \texttt{DENIED}. Полный протокол реального
запуска программы приведён ниже (см. \autoref{lst:mongo-output}).
Администратор проходит любую проверку; \texttt{labuser} может читать и
писать в \texttt{labdb}, но не может ни создать нового пользователя, ни
обратиться к посторонней базе \texttt{otherdb}; \texttt{labguest} может
прочитать только \texttt{public\_notes} -- чтение \texttt{private\_notes}
и любая запись отклоняются. Отдельно проверено, что клиент без учётных
данных не может выполнить вообще ни одной операции -- \texttt{--auth}
действительно принудительно применяется сервером, а не является лишь
опцией по умолчанию.

\codelistingfile[fontsize=\scriptsize]{text}{lab/code/mongo_output.txt}{Реальный протокол работы программы Lab1.2-MongoDB}{mongo-output}

\subsection*{4. Удаление}

Завершающий этап удаляет контейнер, именованный том с данными и сам
образ \texttt{mongo:7} -- после него от установки MongoDB не остаётся ни
одного файла и ни одного объекта Docker.

\labpartbreak{Реализация: OpenFire}

Проект \texttt{Lab1.2-OpenFire} устроен по той же схеме
(\texttt{Install} $\to$ \texttt{Configure} $\to$ \texttt{Verify} $\to$
\texttt{Cleanup}), но, в отличие от MongoDB, официального Docker-образа
у OpenFire нет, поэтому установка выполняется «с нуля» -- пакетом
\texttt{.deb} внутри обычного контейнера \texttt{ubuntu:24.04} -- в точности
так, как это делал бы администратор на настоящей виртуальной машине.

\subsection*{1. Установка}

В контейнер устанавливается пакет \texttt{default-jre-headless} (Java,
необходимая OpenFire), после чего с официального сайта
\texttt{download.igniterealtime.org} загружается пакет
\texttt{openfire\_5.1.2\_all.deb} и устанавливается командой
\texttt{dpkg~-i}. Туда же устанавливается виртуальное окружение Python с
библиотекой \texttt{slixmpp} -- она используется не для настройки
OpenFire, а исключительно для последующей проверки прав доступа
непосредственно по протоколу XMPP (см. пункт 3).

\subsection*{2. Настройка контроля доступа}

OpenFire обычно настраивается через мастер первого запуска в
веб-интерфейсе, что не поддаётся автоматизации без эмуляции браузера.
Вместо этого используется официально документированный, но
малоизвестный механизм \texttt{<autosetup>} в конфигурационном файле
\texttt{openfire.xml}: если в нём присутствует блок \texttt{autosetup}
с \texttt{run=true}, сервер проходит мастер настройки автоматически при
первом запуске и удаляет этот блок из файла. Таким способом заранее
задаются пароль администратора \texttt{admin} и создаётся обычная учётная
запись \texttt{iuser}, а в конфигурацию явно прописывается анонимная
XMPP-аутентификация -- она нужна учётной записи «гость», у которой нет
и не должно быть пароля.

После запуска сервера программа подключается к нему уже не через HTTP, а
непосредственно по протоколу XMPP (библиотека \texttt{slixmpp}, роль
администратора): создаёт комнату \texttt{labroom}, включает в её
настройках модерацию (\texttt{moderatedroom}) и явно назначает
\texttt{iuser} аффилиацию \texttt{member} -- то есть именно администратор
предоставляет пользователю право голоса в этой комнате
(см. \autoref{lst:openfire-muc-setup}). Учётная запись «гость»
отдельно не создаётся вовсе: анонимный участник входит в ту же комнату
без пароля и получает роль «посетитель», которая в модерируемой комнате
не даёт права говорить -- только читать.

\begin{codelisting}[fontsize=\footnotesize]{python}{Настройка модерируемой комнаты и назначение пользователю права голоса}{openfire-muc-setup}
form = await self.plugin["xep_0045"].get_room_config(ROOM)
form.set_values({
    "muc#roomconfig_roomname": "Lab Room",
    "muc#roomconfig_moderatedroom": True,
    "muc#roomconfig_publicroom": True,
    "muc#roomconfig_persistentroom": True,
    "muc#roomconfig_membersonly": False,
    "muc#roomconfig_whois": "moderators",
})
await self.plugin["xep_0045"].set_room_config(ROOM, form)

# Admin grants the "user" identity full (voice) rights in this subsystem.
await self.plugin["xep_0045"].set_affiliation(ROOM, affiliation="member", jid=USER_JID)
\end{codelisting}

\subsection*{3. Проверка корректности}

Проверка охватывает обе подсистемы контроля доступа OpenFire независимо
друг от друга. Доступ к веб-консоли администрирования проверяется
реальным HTTP-логином через \texttt{curl} (с извлечением CSRF-токена из
формы и последующей отправкой формы входа): \texttt{admin} успешно
получает сессию и перенаправление на \texttt{index.jsp}, а \texttt{iuser}
с верным паролем получает страницу с ошибкой «\texttt{Login failed... make
sure... you're an admin}» -- OpenFire различает «валидная учётная запись»
и «учётная запись администратора» на уровне самой веб-консоли. Доступ к
чат-комнате проверяется на уровне протокола: небольшой XMPP-клиент на
Python подключается под нужной личностью, входит в комнату
\texttt{labroom} и пытается отправить сообщение, а затем проверяет, не
вернул ли сервер ответной ошибки \texttt{forbidden}. Полный протокол
реального запуска приведён ниже
(см. \autoref{lst:openfire-output}): и \texttt{admin}, и \texttt{iuser}
могут говорить в комнате, тогда как анонимный гость и подключение с
несуществующей учётной записью получают отказ.

\codelistingfile[fontsize=\footnotesize]{text}{lab/code/openfire_output.txt}{Реальный протокол работы программы Lab1.2-OpenFire}{openfire-output}

\subsection*{4. Удаление}

Поскольку OpenFire был установлен не в отдельный именованный том, а
непосредственно в файловую систему контейнера, для полного удаления
достаточно одной команды \texttt{docker rm~-f} -- вместе с контейнером
удаляются и сам сервер, и встроенная база данных, и созданные плагины,
без необходимости отдельно подчищать файлы на диске.
